\documentclass[11pt,aspectratio=169]{beamer}

% Русский язык. Этот вариант преамбулы рассчитан на pdfLaTeX.
\usepackage[T2A]{fontenc}
\usepackage[utf8]{inputenc}
\usepackage[english,russian]{babel}

% Математика, рисунки и таблицы.
\usepackage{amsmath,amssymb}
\usepackage{graphicx}
\usepackage{booktabs}
\usepackage{stmaryrd}

% Стандартное оформление Beamer.
% Для другой встроенной темы замените default, например, на Madrid.
% Для формата 4:3 замените aspectratio=169 на aspectratio=43 выше.
\usetheme{Madrid}
\usecolortheme{default}
\usefonttheme{default}
\setbeamertemplate{navigation symbols}{}
\setbeamertemplate{footline}[frame number]

% Данные титульного слайда: замените на свои.
\title[Краткое название]{Формирование заказов торговой сети с использованием специализированных алгоритмов прогнозирования с экспертами}
\subtitle{}
\author[И.~О.~Фамилия]{Е.~А.~Цыбулин}
\institute{МГУ имени М.~В.~Ломоносова}
\date{\today} % Можно указать дату явно или оставить \date{}.

% При необходимости раскомментируйте оглавление перед каждым разделом.
% \AtBeginSection[]{
%   \begin{frame}{План раздела}
%     \tableofcontents[currentsection]
%   \end{frame}
% }

\begin{document}

% Титульный слайд не входит в нумерацию.
\begin{frame}[plain,noframenumbering]
  \titlepage
\end{frame}

\begin{frame}{План презентации}
  \tableofcontents
\end{frame}

\section{Введение}

\begin{frame}{Базовые определения}
  \medskip
  \begin{block}{Определение}
    \begin{itemize}
      \item \textit{Временной ряд} — последовательность некоторых событий во времени. Множество исходов — $\mathcal{Y}$.
      \item \textit{Эксперты} в данной работе — чёрный ящик, который выдаёт предсказания следующего события из множества допустимых прогнозов $\mathcal{D}$.
    \end{itemize}
  \end{block}

  \begin{itemize}
    \item Природа значений временного ряда необязательно имеет вероятностный характер. Последовательность может быть случайной, детерминированной или враждебной по отношению к прогнозирующему алгоритму.
  \end{itemize}
  
  \begin{exampleblock}{Пример}
    Прогноз потребления электроэнергии или формирование заказов для поставщика товаров.
  \end{exampleblock}
  \medskip
\end{frame}

\section{Основная часть}

%\begin{frame}{Постановка задачи}
%  Имеется множество исходов $\mathcal{Y}$ и множество допустимых прогнозов $\mathcal{D}$.
%  \begin{block}{Протокол прогнозирования}
%    \begin{enumerate}
%      \item Среда имеет историю значений $y_1, \dots, y_{t-1}$ из $\mathcal{Y}$;
%      \item Эксперты $\{1, \dots, N\}$ делают прогнозы $\gamma_t^{(1)}, \dots, \gamma_t^{(N)}$ из $\mathcal{D}$;
%      \item Смешивающий алгоритм (далее распределитель) изучает решения экспертов и предлагает свой прогноз $\gamma_t \in \mathcal{D}$;
%      \item Среда раскрывает следующее значение $y_t \in \mathcal{Y}$;
%      \item Эксперты и распределитель получают свои потери $l_t^{(1)}, \dots, l_t^{(N)}$.
%    \end{enumerate}
%  \end{block}
%\end{frame}

\begin{frame}{Постановка задачи}
  На каждом шаге нам известны только потери экспертов (знать функцию потерь необязательно).

  В моменты времени $t = 1, 2, \dots$:
  \begin{itemize}
    \item Распределитель находит вектор распределения экспертов $\bar{p_t} = (p_t^1, \dots, p_t^N)$, где $\sum\limits_{i=1}^{N}p_t^i = 1$, $p_t^i \geqslant 0$.
    \item Эксперты объявляют потери $\ell_t^i$;
    \item Распределитель подсчитывает свои потери $\ell_t = (\bar{p_t} \cdot \bar{\ell_t}) := \sum\limits_{i=1}^{N}p_t^i \cdot \ell_t^i$;
    \item Распределитель пересчитывает веса экспертов.
  \end{itemize}
  
\end{frame}

\begin{frame}{Функции потерь}
  \begin{block}{Определение: функция потерь}
    Функция потерь — некоторая функция
    \[\ell: \mathcal{D} \times \mathcal{Y} \to [0, +\infty),\]
    выпуклая по множеству $\mathcal{Y}$.
  \end{block}

  \begin{exampleblock}{Примеры}
    \[\ell(\gamma, y) = \lambda_1 |\gamma - y|_+ + \lambda_2 |y - \gamma|_+,\]
    где $|x|_+ := \max{(x, 0)}$;
    \[\ell(\gamma, y) = (\gamma - y)^2, \quad \ell(\gamma, y) = \llbracket \gamma \neq y \rrbracket.\]
  \end{exampleblock}
\end{frame}

\begin{frame}{Оценка алгоритма}
  \begin{block}{Определение: кумулятивная потеря}
    \[L_T = \sum\limits_{t=1}^{T} \ell(\gamma_t, y_t), \quad L_T^i = \sum\limits_{t=1}^{T} \ell(\gamma_t^i, y_t)\]
  \end{block}

  \begin{block}{Определение: регрет относительно лучшего фиксированного эксперта}
    \[R_T = L_T - \min\limits_{1 \leqslant i \leqslant N} L_T^i\]
  \end{block}

  \begin{alertblock}{Замечание}
    Лучший фиксированный эксперт выбирается по суммарной потере за все $T$ шагов;
    его индекс не меняется во времени.
  \end{alertblock}
\end{frame}

\begin{frame}{Hedge}
  Начальное распределение весов: $w_1^i = \frac{1}{N}$. Для моментов времени $t = 1, 2, \dots$:
  \begin{itemize}
    \item Распределитель вычисляет $\bar{p_t} = \frac{\bar{w_t}}{\sum\limits_{i=1}^{N}w_t^i}$;
    \item Эксперты объявляют свои потери $\ell_t^i$. Пусть $\bar{\ell_t} = (\ell_t^1, \dots, \ell_t^N)$;
    \item Распределитель подсчитывает свои потери $\ell_t = (\bar{p_t} \cdot \bar{\ell_t}) := \sum\limits_{i=1}^{N}p_t^i \cdot \ell_t^i$;
    \item Распределитель пересчитывает веса экспертов $$w_{t+1}^i = w_t^i e^{-\eta \ell_t^i},$$ где $\eta > 0$ — параметр обучения.
  \end{itemize}
  Кумулятивные потери:
  \[H_t = \sum\limits_{s=1}^{t} \ell_s, \quad L_t^i = \sum\limits_{s=1}^{t} \ell_t^i.\]
\end{frame}


\begin{frame}{Две колонки и таблица}
  \begin{columns}[T,onlytextwidth]
    \begin{column}{0.47\textwidth}
      \textbf{Порядок работы}
      \medskip
      \begin{enumerate}
        \item Задать исходные данные.
        \item Применить метод.
        \item Оценить результат.
      \end{enumerate}
      % Для появления пунктов по очереди используйте
      % \begin{enumerate}[<+->] вместо \begin{enumerate}.
    \end{column}

    \begin{column}{0.47\textwidth}
      \textbf{Обозначения}
      \medskip

      \begin{tabular}{@{}ll@{}}
        \toprule
        Символ & Значение \\
        \midrule
        $n$       & Число наблюдений \\
        $x_i$     & Наблюдение $i$ \\
        $\bar{x}$ & Среднее значение \\
        \bottomrule
      \end{tabular}
    \end{column}
  \end{columns}
\end{frame}

\begin{frame}{Иллюстрация}
  \begin{figure}
    \centering
    % Положите рисунок my-figure.pdf рядом с main.tex.
    % Можно заменить имя файла на своё (PDF, PNG или JPG).
    \IfFileExists{my-figure.pdf}{%
      \includegraphics[width=0.75\textwidth,height=0.52\textheight,
                       keepaspectratio]{my-figure.pdf}%
    }{%
      \fbox{%
        \parbox[c][0.44\textheight][c]{0.72\textwidth}{%
          \centering Место для рисунка%
        }%
      }%
    }
    \caption{Краткая подпись к иллюстрации}
  \end{figure}
\end{frame}

\section{Заключение}

\begin{frame}{Заключение}
  \begin{itemize}
    \item Основной результат работы.
    \item Область применения и ограничения.
    \item Направление дальнейшего исследования.
  \end{itemize}

  \bigskip
  \centering
  Спасибо за внимание!
\end{frame}

% Заготовка нового слайда: скопируйте и раскомментируйте.
% \begin{frame}{Заголовок слайда}
%   Содержимое слайда.
% \end{frame}

\end{document}
